\begin{table*}[t]
\centering
\caption{Proposed tiered release-testing protocol for nanocomposite and photocatalytic marine coatings. The protocol is a methodological recommendation synthesized from the reviewed release, cleaning, and coating studies; it is not a reported dataset.}
\label{tab:release_protocol}
\small
\begin{tabularx}{\textwidth}{l l X X X}
\toprule
Tier & Exposure design & Fractions and measurements & Required controls & Decision output \\
\midrule
1: baseline immersion & Closed-bottle natural or artificial seawater; report exposed area, coating mass, water-to-area ratio, salinity, pH, temperature, dissolved oxygen, light, and sampling times & Dissolved, colloidal, and particulate fractions; ICP-MS/OES for elements; particle number and size by single-particle ICP-MS, microscopy, or validated orthogonal method; coating mass balance & Field and procedural blanks, unfilled matrix, neat substrate, duplicate samples, matrix spikes, filtration recovery, and particle-loss check & Area-normalized concentration--time profile and fraction-specific mass balance \\
2: service stress & Add representative flow, calibrated abrasion or waterjet exposure, wet--dry cycling, UV, and repeated cleaning; document stress energy and exposure history & Repeat Tier 1 measurements plus mass loss, roughness, wettability, adhesion, coating integrity, and antifouling endpoint before and after stress & Untreated coating, stress-only blank, positive release control, and pre/post surface characterization & Retained performance versus release and wear under a defined service scenario \\
3: mesocosm or field & Matched panels or coupons in a controlled mesocosm or field deployment; recover water, sediment or biofilm, and detached solids & Dissolved and particulate release, organism/community response, fouling, surface condition, and environmental chemistry; record hydrodynamic and cleaning history & Spatial and temporal replicates, reference panels, recovery efficiency, chain of custody, and receiving-water background & Asset-relevant exposure estimate with uncertainty and ecological context \\
Ecotoxicity module & Use receiving-environment-relevant primary producer or algae, invertebrate or larval, and microbial endpoints; separate dissolved-ion and particle-only exposures where feasible & Concentration--response and time-response data for matched fractions; report measured rather than nominal exposure where possible & Solvent, seawater, dissolved-ion, particle-only, and positive toxicity controls; blinded sample coding where practical & Hazard characterization linked to the same fractions and concentrations measured in release testing \\
\bottomrule
\end{tabularx}
\end{table*}
